\paragraph{C.3z (Proved): removing the algebraizability restriction for the ordinary degree-one theorem.}
Let $k$ be an algebraically closed field and let $(R,\mathfrak m,k)$
be an excellent henselian normal noetherian local $k$-algebra whose
specified coefficient map induces the identity on the residue field.
Assume that $X=\operatorname{Spec}R\setminus\{\mathfrak m\}$ is regular.
There is a normal local $k$-algebra $A_0$ essentially of finite type,
with residue field $k$, and a local regular $k$-homomorphism
\[
                    A_0\longrightarrow R
\]
which induces an isomorphism of maximal-ideal completions.
Consequently $\operatorname{Spec}R$ admits a proper generically finite
surjective morphism from a regular integral scheme.

If $N=\dim R\ge3$, the point-realization conclusion of C.3w holds
for this $R$: every $k$-point of every closed abelian subvariety
of $\mathbf{Pic}^{\mathrm{loc}}_{R/k}$ belongs to
$\operatorname{Pic}(X)_{\mathrm{alg}}$ for ordinary algebraic
equivalence on the fixed product.
In positive characteristic $p$, if $N\ge4$, $0\ne f\in\mathfrak m$,
and $Y=X\cap V(f)$ is regular, then
\[
 \operatorname{Pic}(X)_{\mathrm{alg}}
       \longrightarrow\operatorname{Pic}(Y)_{\mathrm{alg}}
\]
is surjective with $p$-primary torsion kernel, and hence is an
isomorphism after tensoring with $\mathbb Q$.
Thus the entire ordinary degree-one table in C.3x extends to
these possibly nonalgebraizable excellent henselian germs.
No statement about mixed characteristic or about higher Chow
groups is included in this extension.

\begin{proof}
\emph{1. Algebraizing the completion with the given coefficient field.}
If $R=k$, we take $A_0=k$ and the identity alteration.
We may therefore assume that the embedding dimension below is positive.
The map $R\to\widehat R$ is regular by excellence.
Normality ascends along this map, so the complete local ring
$\widehat R$ is a normal domain of dimension $N$.
Its puncture is regular: a nonmaximal prime of $\widehat R$
cannot lie over $\mathfrak m$, because
$\widehat R/\mathfrak m\widehat R=k$; at its contraction
the source ring is regular, and regularity ascends along
a regular ring map.

Choose finitely many generators of $\mathfrak m/\mathfrak m^2$.
Their lifts and the specified copy of $k$ give a presentation
\[
             \widehat R\simeq k[[x_1,\ldots,x_e]]/I.
\]
Surjectivity follows by successively lifting homogeneous
polynomials in the generators and using completeness.
Since $k$ is perfect, the formal-power-series Jacobian
criterion identifies regularity of each localization of
this quotient with the required full-rank Jacobian minor.
The exact version we use is
\citet[Proposition 2.5, with the definition of $\mathcal B$
in Section 2]{rootHauserWoblistin2021}: take
$\mathcal B=k[[\mathrm t,\mathrm y]]$, simply relabeling
the $e$ variables. Thus the presentation is formally
smooth away from its closed point in the precise sense
of \citet[Theorem 3.8, p.~32]{rootArtin1969}.
That theorem supplies a finite-type $k$-scheme $V$, a
$k$-rational point $v$, and a $k$-isomorphism
\begin{equation}\label{eq:C3z-formal-algebraization}
                  \widehat{\mathcal O}_{V,v}\simeq\widehat R.
\end{equation}
The local ring $A_0=\mathcal O_{V,v}$ is normal by
faithfully flat descent from its normal completion.
After shrinking $V$ around $v$ we may and do take $V$
normal and integral. Here we use openness of the normal
locus for a finite-type algebra over a field; no
regularity assertion away from $v$ is needed for the
alteration below.

\emph{2. Obtaining an actual map to the original henselian ring.}
Write an affine neighborhood of $v$ in the form
$\operatorname{Spec}(k[t_1,\ldots,t_a]/(F_1,\ldots,F_b))$,
translating the coordinates so that $v$ is the origin.
The formal isomorphism \eqref{eq:C3z-formal-algebraization}
gives a solution $(\widehat t_1,\ldots,\widehat t_a)$
of these polynomial equations in $\widehat R$.
The general henselian G-ring approximation theorem
\citep[Theorem 16.13.1, Tag 07QY]{rootStacksRegularApproximation}
provides a solution $(t'_1,\ldots,t'_a)$ in $R$ with
\[
                   t'_j-\widehat t_j\in\mathfrak m^2\widehat R.
\]
This gives an actual local $k$-map $A_0\to R$.
It extends to a continuous map
$\widehat A_0\to\widehat R$.
The coordinate images generate the cotangent space
of $\widehat R$, since they agree modulo
$\mathfrak m^2$ with those of the given isomorphism.
The same successive approximation by polynomials now
shows that $\widehat A_0\to\widehat R$ is surjective.
Its kernel is prime, both rings are domains of the
same finite dimension $N$, and a nonzero prime in a
domain strictly lowers the maximal possible chain
length in its quotient. The kernel is therefore zero.
We have constructed a map inducing a genuine
isomorphism of completions, not merely matching
some chosen finite jet.

The composite $A_0\to R\to\widehat R$ is regular:
under this new isomorphism it is the completion map
of the excellent ring $A_0$.
The second arrow $R\to\widehat R$ is faithfully flat.
It follows that $A_0\to R$ is regular, by
\citet[Lemma 15.42.7, Tag 07NT]{rootStacksRegularApproximation}.
For clarity, that descent first gives flatness;
on each fibre, faithful flatness descends geometric
regularity. This proves the first assertion.

\emph{3. A regular proper alteration over $R$.}
Apply \citet[Theorem 4.1 and Remark 4.2, p.~66]{D_DeJong1996Alterations}
to the normal algebraic neighborhood $V$.
It gives a regular integral scheme $Z_0$ and a proper
generically finite surjective map $Z_0\to V$.
Base change first to $A_0$ and then by the actual
regular map $A_0\to R$.
Regular maps remain regular under finite-type base
change \citep[Lemma 15.42.3, Tag 07BY]{rootStacksRegularApproximation},
so the resulting scheme $Z_R$ is regular.
It is noetherian and proper over $R$; its generic
fibre is finite and nonempty because $A_0\to R$
is faithfully flat and the original alteration is
generically finite. A regular noetherian scheme
has disjoint integral connected components.
Choose one component meeting this generic fibre.
It is proper and generically finite over the
integral local base; its image is closed and
contains the generic point, hence it is surjective.
This is the claimed regular integral alteration.
We neither assert regular alterations for all
excellent schemes nor use altered local uniformization
as if it supplied a single global regular alteration.

\emph{4. Transfer to ordinary algebraic equivalence.}
The only use of an algebraic model in C.3w is its
step 1, which supplies one regular integral proper
alteration. Step 3 above supplies exactly that input
over the present $R$. Every subsequent hypothesis
is unchanged: the Stein algebra is a finite normal
local algebra with residue $k$; tensor products with
smooth parameter curves are noetherian and regular;
proper formal existence applies over $\widehat R$;
and the whole-family approximation in C.3w uses the
same general theorem 07QY, already valid for this
excellent henselian $R$.
The exceptional-divisor lattice, formal-unit gerbe,
proper-curve rigidity and finite-flat norm arguments
of steps 2--7 therefore prove the stated
point-realization assertion with this replacement of step 1.

\emph{5. Lifting target abelian subgroups without algebraizing the divisor.}
We supply the two replacements needed to use the argument of C.3q;
we do not assume that $f$ descends through $A_0\to R$.
Put $G=\mathbf{Pic}^{\mathrm{loc}}_{R/k}$ and let
$H=\prod_a\mathbf{Pic}^{\mathrm{loc}}_{C_a/k}$, where
$\prod_a C_a$ is the finite normalization of $R/(f)$.
These are locally finite-type group schemes in dimension $N\ge4$
by the precise Boutot input in C.3n. Restriction gives $r:G\to H$.

First, the functor comparison proved separately in C.3p gives
\[
 G\simeq\mathbf{Pic}^{\mathrm{loc}}_{\widehat R/k}
   \simeq\mathbf{Pic}^{\mathrm{loc}}_{A_0^h/k}.
\]
Here this is a comparison of group schemes, including identity
components, rather than only a bijection on $k$-points.
Indeed, for every finite-type $k$-algebra $D$, the henselized
tensor pairs for $R$ and $A_0^h$ have the same ideal-adic completion
\[
             \varprojlim_n(R/\mathfrak m^n)\otimes_k D.
\]
They are noetherian henselian pairs, so the Picard corollary of
\citet[Section III, p.~573]{D_Elkik1973} identifies the Picard
groups of their ideal complements with that of this common
completed complement. The maps are natural in $D$ and survive
etale sheafification. The affine finite-type open subschemes
of the represented group schemes then test the inverse by
Yoneda. The given coefficient maps and the $k$-isomorphism of
completions ensure that this comparison is over the specified $k$.
Thus the finite generation of the local Neron--Severi point
group for the normal algebraic model
\citep[Definition 3.4 and the following paragraph, p.~468]
{D_Kollar2016LocalPicard} gives the same finite generation
for $G/G^0$. Only bounded prime-to-$p$ torsion is needed.

Choose a prime $\ell\ne p$ avoiding that finite torsion group.
For any closed abelian subvariety $B\subset H$, C.3p lifts
every $B[\ell^n](k)$ uniquely into $G(k)[\ell^n]$.
Its component class is zero, so the lift lies in
$G_0=(G^0)_{\mathrm{red}}$. The closed image $r(G_0)$
contains the Zariski-dense union of these torsion groups;
therefore it contains $B$. As in C.3q, taking the reduced
identity component of $r^{-1}(B)\cap G_0$ gives a smooth
connected commutative algebraic group $E$ and a faithfully
flat homomorphism $r_E:E\to B$.

Second, its scheme-theoretic kernel $K$ is of finite type
and C.3k(i) gives $K(k)[\ell']=0$ for every prime
$\ell'\ne p$. We now prove that some $[p^e]$ kills $K$
as a group scheme. This assertion allows an infinitesimal
multiplicative part such as $\mu_p$.
The reduced identity component $U=(K_{\mathrm{red}})^0$
is unipotent. To see this, apply the smooth connected
Chevalley decomposition to $U$. Its connected affine
part is an extension of a torus by a unipotent group.
A nonzero torus already supplies prime-to-$p$ torsion.
If the abelian quotient were positive dimensional,
its prime-to-$p$ torsion would lift to torsion in $U(k)$,
because multiplication by that prime is surjective on
the connected affine part. Both possibilities contradict
the displayed torsion exclusion. This is the argument
of \citet[Lemma 2.6 and proof, p.~466]{D_Kollar2016LocalPicard},
applied to the reduced group.
The same lifting observation, using the invertibility
of prime-to-$p$ multiplication on $U$, shows that the
finite component group of $K_{\mathrm{red}}$ is a $p$-group.
The group $U$ embeds in an upper unitriangular matrix
group and is killed by a power of $p$. Hence a power
$[p^a]$ kills all of $K_{\mathrm{red}}$.

The reduced group is affine: it is a finite disjoint
union of translates of $U$. Its nilpotent thickening
$K$ is affine as well
\citep[Lemma 37.2.3, Tag 06AD]{rootStacksRegularApproximation}.
Choose a faithful closed representation $K\hookrightarrow
\operatorname{GL}_d$
\citep[Theorem 4.9 and Corollary 4.10, p.~87]
{rootMilneAlgebraicGroups2022}.
In the universal matrix of $[p^a]$ on $K$, all entries
of $M=[p^a]-I_d$ lie in the nilradical $J$ of $\mathcal O(K)$,
since this map is the identity matrix on $K_{\mathrm{red}}$.
Choose $s$ with $J^s=0$ and $b$ with $p^b\ge s$.
In characteristic $p$ we have
\[
                    (I_d+M)^{p^b}=I_d+M^{p^b}=I_d.
\]
Faithfulness proves $[p^{a+b}]_K=0$ scheme-theoretically.

Write $e=a+b$. The homomorphism $[p^e]_E$ descends
through the faithfully flat quotient $E\to B$ to a
homomorphism $\widetilde s:B\to E$ with
$r_E\widetilde s=[p^e]_B$. Its image $A\subset E\subset G$
is a closed abelian subvariety; the maps $B\to A\to B$
are isogenies with $p$-primary kernels, exactly as in
the final descent argument of C.3q.
Every ordinary algebraically trivial class of $Y$
belongs to some such $B(k)$ by the actual abelian-family
lemma of C.3n, applied to its finitely many regular
components. The surjection $A(k)\to B(k)$ and step 4
realize a preimage by an ordinary algebraic equivalence
on $X$. This proves integral surjectivity without any
simultaneous algebraization of $(R,f)$.

C.3n already has the general coefficient-field
scope used here and gives the $p$-primary kernel.
The counterexamples in dimensions two and three in
C.3x lie in this larger class; the dimension-one
vanishing and the $\overline{\mathbb F}_p$ torsion
argument also retain their original general scope.
This proves the extension of the entire table.
\end{proof}
